% Appendix: thermal (IR) imaging study, 4 Nov 2026 onward. Draft 27 Nov 2026 (writer id: thermal).
% Every number and dated fact: chapters/app_thermal_NOTES.md (source path and line).
% Points to Chapter 10, Section "What thermal imaging adds" (sec:deploy-thermal) for the A3 drone-thermal result; does not repeat it.
\chapter{Thermal imaging study}\label{app:thermal}

This appendix records the thermal (infrared, IR) imaging work from 4 November 2026 on: a one-day desk study, the public data
sets that pair EL and IR images of the same modules, the ranked ideas and the three picks, the readouts of picks 1 and 3 on
public data, and the later use of thermal in the severity track. The thesis reports the one thermal result that it uses, the
module-level value of drone thermal on the DTU set A3, in Section~\ref{sec:deploy-thermal}; this appendix points to it and does not
repeat it. All work here is from November 2026, after the benchmark protocol era began, so no old-era number appears. No IoU is
reported. No model was trained on thermal data, nothing was run on a GPU, and no image of the legacy test set or of E1 was opened.

\paragraph{Conventions.} Every verdict below is the one fixed in a frozen plan before its numbers were computed. A ``robustness
check'' is printed beside a verdict and decides nothing; a ``reading'' is interpretation and is marked as ours. Several rules were
changed after the data had been seen; each such change is marked ``post hoc''. The study used these measures, all at cell level unless
stated:
\begin{itemize}
\item $\Delta T$: a cell's median temperature minus the module median of cell medians, in $^{\circ}$C.
\item EL dark fraction: the share of a cell's inner pixels below 0.5 times the module median, after a small morphological opening.
A cell with a dark fraction of 5\% or more is ``EL-dark''. This is a plain EL threshold, used in place of our model's inactive
prediction because model inference on these sets was on hold (Section~\ref{sec:thermal-decisions}).
\item AUROC (cooler): how well $-\Delta T$ ranks EL-dark cells above the others (0.5 is chance). $\rho$: Spearman rank correlation of
$\Delta T$ with the dark fraction.
\end{itemize}

\section{4 November, morning: the brief and the desk study}\label{sec:thermal-desk}

On 4 November a brief from a collaborator, dated 3 November, asked three questions: (1) are there public data sets that pair EL with thermal images (also
EL with photoluminescence (PL), UV fluorescence (UVF) or visible images) of the same PV modules; (2) in what ways could thermal help
our EL segmentation, including the industrial pattern ``drone IR for screening, EL for detail''; and (3) could the industry partner
or IIT Mandi capture paired EL and thermal images of a few modules. The brief asked for an optimistic, generative search, with every
claim tagged as measured in our files, taken from a source, or judgement. The rules were strict: a desk study first (web pages,
dataset metadata and abstracts only, no downloads), no GPU, no outside contact, no reads of test or E1 data, and no request for
hand labels. The work was split into roles: a data search, an idea list, a capture plan, a synthesis that ranked the ideas, and an
independent checker that verified each cited fact. The team started at 11:56.

\paragraph{What thermal can and cannot show.} EL shows where injected current reaches the silicon and recombines; thermal shows where
power turns into heat, blurred by heat spreading. Table~\ref{tab:thermal-physics} gives the mapping the study wrote for our classes.
A crack that isolates nothing (mode A in the scheme of K\"ontges et al.~\citep{kontges2011risk}; the mapping is ours) is expected to
be thermally silent: Dhimish and Lazaridis found no hot spot on such cells~\citep{dhimish2021hotspots}, and in drone IR of
lab-damaged modules cracked modules were ``difficult to distinguish from healthy ones'', at about 1\,$^{\circ}$C or less
\citep{delpradosantamaria2025drone}. An area isolated by cracks should read cooler than its neighbours under forward bias in the
dark (no current reaches it) and warmer in sunlight (it cannot export its power); both are device-physics expectations, not
measurements. A study of aged modules from the European Commission's Joint Research Centre found that dark IR gives ``few local
information'' while illuminated IR separates damage classes in line with EL \citep{berardone2017jrc},
so a weak result in dark IR was not to be read as a verdict on illuminated IR. Drone IR at 3\,cm per pixel or finer resolves about 5
to 16 or more pixels across a 156\,mm cell, a cell-level signal at best.

\begin{table}[htbp]
\centering\small
\caption{What each modality is expected to show for our classes, as written in the desk study of 4 November 2026. ``Expected''
entries are device-physics judgement; entries with a citation are from that source.}
\label{tab:thermal-physics}
\begin{tabular}{@{}p{3.2cm}p{2.3cm}p{3.0cm}p{3.6cm}@{}}
\toprule
Class & EL & Dark forward-bias IR & Sunlit (drone) IR \\
\midrule
Crack that isolates nothing & thin dark line & none expected & no hot spot~\citep{dhimish2021hotspots}; at most about 1\,$^{\circ}$C~\citep{delpradosantamaria2025drone} \\
Inactive area (isolated by cracks) & dark region & cooler region expected & warmer region expected \\
Scuff (not one of our classes) & long dark streak & unknown; silent if surface-only & unknown \\
Shunt, potential-induced degradation & dark cell or spots & hot spot & warm cell \\
Interconnect disconnection & dark region & & visible from 600~W/m$^2$~\citep{delpradosantamaria2025drone} \\
\bottomrule
\end{tabular}
\end{table}

The net reading of the desk study (our reading, as the study wrote it): thermal is a cell-level signal for inactive area, hot spots
and severity, and a weak or absent signal for crack lines and scuff; it sees the electrical consequence, not the crack. It could
therefore help as a training-time teacher, a weak label, a triage stage or a severity output, but not as a crack-pixel input. Since
neither BenchmarkEL, the legacy test set nor E1 has thermal images, any gain on our benchmarks would have to come from an EL-only
model trained with paired data. The same reasoning meant thermal could not fix the E1 gap directly: the team's brief named that gap as
missed cracks in half-cut cells with scuff streaks (the hard stratum of Chapter~\ref{ch:t3b}), and both thin cracks and, probably,
scuff are thermally silent.

\section{4 November, 11:57 to 12:53: the public paired data sets}\label{sec:thermal-data}

The data search found no public set with pixel-registered EL and thermal images at cell level and crack or inactive masks. Three
open sets, all under CC~BY~4.0, hold EL and IR images of the same physical modules (Table~\ref{tab:thermal-sets}); the pairing is by
module, so cells must be aligned by our own registration. Together they hold about 215 modules with IR, about 203 with both EL and
IR. The table also listed 5 EL plus PL rows, 4 EL plus visible or near-IR rows (only one of which, the staged-damage set of Polenghi
et al.~\citep{polenghi2026data}, pairs the same cell), 12 thermal-only sets with defect labels, and 12 sets described in papers but not
released.

\begin{table}[htbp]
\centering\small
\caption{The three public sets with EL and IR of the same modules (4 November 2026). Sizes are of the full public folders; what was
downloaded is in Section~\ref{sec:thermal-decisions}.}
\label{tab:thermal-sets}
\begin{tabular}{@{}p{1.9cm}p{3.6cm}p{3.3cm}p{4.2cm}@{}}
\toprule
Set & Modules & IR & EL and extras \\
\midrule
A1 (OSF csr6t; \citet{colvin2023degradation}) &
10-year rooftop, mono Al-BSF, 36 cells in 3 $\times$ 12; 156 characterised, IR for 157 &
dark forward bias at 3.39\,A (about 1 $\times$ Isc), 180\,s by the file names (the paper says 5 min); 240 $\times$ 320 matrices in $^{\circ}$C; 0.26\,GB &
EL at 6 currents, 14.95\,GB; UVF 2.01\,GB; I-V; EL-derived per-cell dark I-V (158 files) \\
A2 (OSF ar76v) & 19-year array, mono, 72 cells in 6 $\times$ 12; 42 modules, IR for 15 (the camera failed) &
dark forward bias at 5.289\,A, 60\,s; 16 captures; 0.02\,GB & EL at 6 currents, 3.82\,GB; UVF; I-V \\
A3 (DTU; \citep{delpradosantamaria2025drone}) & lab-damaged: cracks 8, interconnects 8, potential-induced degradation 6,
bypass diode 7, glass 2, soiling 6 + 6; 43 in all &
drone, at 8, 10, 14 and 20\,m and 200 to over 800~W/m$^2$; arrays 605 $\times$ 363 in $^{\circ}$C &
one EL image per degraded module; STC I-V (V, I only); archive 642.6\,MB \\
\bottomrule
\end{tabular}
\end{table}

Three facts about these sets shaped everything after. A1 and A2 are old full-cell modules, with no half-cut cells and no scuff. A1's
losses are ``primarily resistive'' according to its authors, so it is a series-resistance set more than an inactive-area set. A3 has
half-cut modules only in its diode, glass and soiling groups, none with a crack. No set is known to show scuff. Our reading at the
time: about 200 paired modules suffice for a feasibility study of cross-modal weak labels and a demonstration pipeline, not for
training a fusion model.

\section{4 November, 12:08 to 12:59: ideas, ranking and the top three}\label{sec:thermal-ideas}

The idea list held 40 ideas (T01 to T40) in 8 families: EL and thermal fusion; cross-modal pretraining and distillation (a thermal
teacher for an EL student); thermal and physics as weak labels; EL to thermal synthesis; thermal for unseen-module generalisation
(including a ``thermal silence'' test for scuff); the industrial pipeline (IR screening, then EL); other modalities (EL against PL
for inactive masks, UVF crack lines against scuff); and data and evaluation (for example a registered A1 cell set as a release).
The list stated its own limits: no direct fix for the scuff and half-cut gap, no gain on test or E1 from thermal inputs, a weak
record of pretraining in this project (Chapter~\ref{ch:legacy}), and cell-level, not pixel-level, signals. After the industry
partner's situation became known (Section~\ref{sec:thermal-decisions}), a second list added 14 drone-only ideas (W01 to W14), such as
an IR-ranked queue of modules for the first EL session with predictions frozen in advance, a module identity map with a half-cut
survey, and a physics-rule anomaly score.

The synthesis merged these with 7 ideas from the data and capture roles into 58 rows and scored each as value $\times$ probability
$/\sqrt{\text{effort}}$; every score is judgement. The top ten were T26 (IR screening, then EL), T37, T15, W03, T27, W02, W01, T02,
X3 and W11. The top three were chosen to spread risk:
\begin{enumerate}
\item \textbf{Pick 1, ``what thermal catches''}: IR screening against EL on A3, A2 and A1, with a shared gate G0 (is there a thermal
signal at cell level?). Public data, CPU only, aimed at the deployment chapter. Kill if the top 20\% of IR-ranked modules hold at
most 20\% of the EL dark area, or if G0 is null in both the dark and the sunlit condition.
\item \textbf{Pick 2, the industry partner's drone thermal now and low-cost EL soon}: a visit day and a night, about INR 27,000 to
30,000 of equipment.
\item \textbf{Pick 3, inactive area from physics labels}: A1's EL-derived per-cell voltage as a label (T15), then thermal as a cell
label (T02, T12), with the registered A1 set (T37) as an output; about 5.5 GPU hours only if gated.
\end{enumerate}

\paragraph{Capture plan.} The first capture plan was one day on 5 to 6 recycler modules (at least two half-cut and two with
scuff-like streaks): EL at 0.1 and 1.0 times Isc, dark IR under the same bias, UVF and visible images, aluminium-tape fiducials for
cell alignment, about INR 45,000 to 55,000 of equipment if an EL rig existed. Its point: an EL-dark streak with no thermal or UVF
signature is a scuff candidate; EL-dark with a thermal anomaly is a crack or inactive candidate. When the partner turned out to have
no EL rig, the plan was redone (Section~\ref{sec:thermal-decisions}): drone IR becomes a triage and pairing tool, not a crack
labeller; and EL could come sooner from a modified low-cost camera with a 300\,W supply (about INR 25,000 plus lens and filter, its
silicon sensitivity near 1,150\,nm unconfirmed). Drone EL would need a short-wave IR camera of about EUR 10,000 to 15,000 and was
ruled out. The plan ended in 20 questions for the partner and IIT Mandi and 3 for dataset authors. None was sent, under the
no-contact rule.

\paragraph{Fact-check.} The checker's first pass, on the data, idea and capture files, covered 110 items: 73 correct, 31 fixed,
6 unverifiable. Most fixes improved the picture; for example A1's IR turned out to be bias-matched forward-bias IR on about 157
modules. The second pass, on the synthesis, covered 33 items: 13 correct, 18 fixed, 2 unverifiable. The top three did not change.

\section{4 November, 12:39 to 13:47: decisions and downloads}\label{sec:thermal-decisions}

At 12:39 a collaborator's decisions arrived. (1) The downloads for picks 1 and 3 were approved, CPU only: A3, an A2 subset
and an A1 subset with the dark I-V files, about 1.2\,GB; a further 0.55\,GB A1 extension still needed approval. (2) The industry
partner has only a drone thermal setup; its EL rig was 2 to 3 months away, so pick 2 had to be re-planned. (3) Three items stayed
open: where our ensemble's inference may run, licences and release, and capture funding. Because of item 3, our model was not run on
these sets in the 4 November study; picks 1 and 3 used the EL dark fraction instead.

The downloads were selected by frozen rules before any content was read. A1: all 159 dark I-V files plus 10 modules drawn by a
seeded rule stratified on a severity score (5 bins $\times$ 2 modules; the most severe module always included), 209 files,
286,886,436 bytes. A2: all IR files plus one EL frame per IR module at the IR current, 47 files, 237,009,189 bytes. A3: the full
archive, 642,610,942 bytes, extracted to 1,414 files. Every file was checked against the publishers' hashes by the data role and,
independently, by the checker. On 12 November, after the laptop's operating system was reinstalled, all three manifests were
checked again with no failure.

\section{4 November, 12:55 to 13:45: pick 1, does dark IR flag EL-dark cells? (G0)}\label{sec:thermal-g0}

\paragraph{Plan.} The plan was frozen at 12:55 to 12:56, before any data value for pick 1 was read, and amended six times, each
amendment frozen before the numbers it governs. Verdict rule (R3): \emph{go} if AUROC (cooler) $\geq 0.8$ with the predicted sign;
\emph{null} if $|\rho| < 0.3$ and AUROC $< 0.7$; anything else is ``between''. The fusion and weak-label families are killed only if
G0 is null in both the dark condition (A1, A2) and the sunlit one (A3). Modules whose registration residual exceeded 5\% of a cell side
were excluded. A regular cell grid looks the same flipped or rotated, so which IR cell is which EL cell needs outside evidence: the
authors' own A1 IR crop is the raw frame rotated by 180$^{\circ}$, so rotation by 180$^{\circ}$ was fixed as primary for A1 and, by
the same lab and rig (our judgement), for A2; the other three orientations were reported beside.

\paragraph{G0 v1 (13:10).} Registration used the team's tool, version 0.2. Results are in Table~\ref{tab:thermal-g0}. A2 was
``between'': the signal had the predicted sign (EL-dark cells cooler) but stayed below the go bar. A1 was null. Pooled was ``between''.
Because the dark condition is not null overall, G0 could not be killed by A3 whatever A3 showed. The checker recomputed every value
from the raw files with its own code and found a match. Robustness checks: on A2, the two vertically flipped orientations carried the
signal and the two unflipped did not; the best of four orientations against a within-module permutation null gave $p = 0.001$ for
both sets. Two weaknesses were noted at the time: the EL images are tone-mapped 8-bit previews of the raw files, so dark fractions are
approximate, and IR cells are only 12 to 15 pixels across.

\paragraph{Post hoc: the signal is in rows and edges.} After the v1 readout, the checker subtracted each module's row medians and then
its column medians from $\Delta T$. The association vanished: AUROC 0.49 to 0.52 on A2 and 0.52 to 0.57 on A1 across orientations.
This was first computed after the readout, so it decides nothing. Our reading, as the study wrote it: the dark-IR signal sits in row
and edge structure, not in single cells; it fits real string or edge degradation but equally fits shared artefacts (convection,
heat sinking at the frame, vignetting), and these data cannot tell them apart.

\paragraph{The instrument problem and G0 v2 (13:43).} At 13:13 the pick 3 role found that the registration residual could pass while
the grid had drifted, on A1 modules whose EL shows dark bands inside cells. The first G0 was therefore marked provisional, and a
``G0 v2, post hoc instrument fix'' was run: the reviewed and frozen tool version 0.4.1, written visual criteria for the grid overlays
fixed before any run, a numeric edge check, and per-module keep or drop calls made from overlays only and co-signed blind by the
checker before any number. The verdict rules did not change. G0 v2 kept 4 of 15 A2 and 4 of 10 A1 modules and gave the same verdicts
as v1 (Table~\ref{tab:thermal-g0}); the checker's recompute matched. With the six borderline modules added (a sensitivity set fixed
before the run), pooled AUROC was 0.706. Row and column centring again removed the A2 signal (AUROC 0.525), so the v1 finding was not a
misregistration artefact. On A1 the centred AUROC (cooler) was 0.343: EL-dark cells read slightly \emph{warmer} than their row and
column peers. Our reading: that is the sign expected for series-resistance defects, which dissipate more heat under forward bias, and it
fits A1's resistive losses; with 4 modules and 36 EL-dark cells it is descriptive only.

\begin{table}[htbp]
\centering\small
\caption{Pick 1, gate G0 on dark forward-bias IR (orientation rotated 180$^{\circ}$). Positives are EL-dark cells. Verdicts by the
frozen rule R3. The row and column centred AUROC is post hoc for this arm and decides nothing; in v1 it was computed by the checker
across orientations (range given).}
\label{tab:thermal-g0}
\begin{adjustbox}{max width=\linewidth}
\begin{tabular}{@{}llrrrrrl@{}}
\toprule
Run & Set & Modules kept & Cells / EL-dark & $\rho$ & AUROC (cooler) & Centred AUROC & Verdict \\
\midrule
v1 & A2 & 9 / 15 & 648 / 230 & $-0.550$ & 0.736 & 0.49 to 0.52 & between \\
v1 & A1 & 7 / 10 & 252 / 50 & $-0.285$ & 0.689 & 0.52 to 0.57 & null \\
v1 & pooled & 16 & 900 / 280 & $-0.451$ & 0.708 & & between \\
v2 & A2 & 4 / 15 & 288 / 75 & $-0.515$ & 0.732 & 0.525 & between \\
v2 & A1 & 4 / 10 & 144 / 36 & $-0.246$ & 0.621 & 0.343 & null \\
v2 & pooled & 8 & 432 / 111 & $-0.403$ & 0.687 & & between \\
\bottomrule
\end{tabular}
\end{adjustbox}
\end{table}

\paragraph{Verdict of pick 1: not a go, not a kill.} The parts of pick 1 on A3 were never run in the 4 November session: the
module-level screening table by defect class (R1), the screening recall against EL (R2, the pick's kill rule) and the sunlit cell-level
arm (R3b). Only A3's structure was discovered, at 13:47 to 13:50; it showed that the module fills the whole IR array, so the outline rule
used for A1 and A2 does not apply, and a sixth amendment set R1 to use the whole array with a 3\% inset.

\section{4 November, 12:46 to 13:43: pick 3, physics labels for inactive area (A1 arm)}\label{sec:thermal-pick3}

\paragraph{T15: is A1's EL-derived cell voltage a usable label?} The plan was frozen before any data value was read (version 1.2 at
12:59, with disclosures; Table~\ref{tab:thermal-pick3}). Per cell, the deficit $d$ is the module median voltage minus the cell's
EL-derived voltage at the highest current (3.392 to 3.393\,A). The frozen kill: coverage below 50\%, or no spread, meaning a robust
standard deviation of $d$ below 2.57\,mV (0.1 thermal voltage) \emph{and} under 1\% of cells with $d$ of one thermal voltage or more.
The second clause was met by construction (the largest module severity, 19.3\,mV, was on record before the read), so the kill rested on
the spread clause alone, and the plan disclosed this before the read. Result at 13:00: 5,688 of 5,688 cells valid, robust standard
deviation 3.26\,mV: \textbf{PASS, narrow}. The checker recomputed it from the raw files and matched. The label D90 ($d \geq 4.23$\,mV)
marks 569 cells in 144 of 158 modules; no cell reaches one thermal voltage. Our reading: this is a diffuse, small resistive deficit,
each module's dimmest tenth of cells at about 1.18 times EL dimming, not inactive area.

\paragraph{T37: can the 10 A1 modules be registered between EL and IR at cell level?} Frozen rule: a module is excluded if its EL or
IR inner residual exceeds 5\% of a cell; the pick is killed if more than 30\% are excluded. Version 1 (13:13) excluded 3 of 10, exactly
30\%: a pass at the line. But the visual check showed one included module misregistered at a passing residual of 4.49\%, and UVF failed
on 9 of 10. A post hoc version 2 (13:28) with the improved tool excluded 4 of 10 by the frozen 5\% rule and failed its own added
grid-to-outline rule on all 10: \textbf{FAIL}. The checker then found the likely cause (the crop geometry is measured; its role in the
failure is our reading): the authors' A1 IR file is a fixed window of 134 $\times$ 35 pixels, not the module, which spans about 145
$\times$ 37 pixels, so the crop trims about 0.9 of a cell, and forcing 12 columns into it gave a pitch error of about 7\%. Any later A1
work must register on the raw IR frame.

A second problem with the label came from the A1 paper: the authors compute each cell's voltage from ``a bright pixel from each cell'',
and ``cells with bright spots in high injection EL show higher RS values''. A degraded cell may thus have a \emph{higher} voltage, so
the frozen ``lower is worse'' direction may be inverted for A1 (our reading).

\begin{table}[htbp]
\centering\small
\caption{Pick 3 on A1 (4 November 2026). Rules frozen before data unless marked post hoc.}
\label{tab:thermal-pick3}
\begin{tabular}{@{}p{2.6cm}p{4.4cm}p{4.4cm}p{2.2cm}@{}}
\toprule
Step & Rule & Result & Verdict \\
\midrule
T15 dark I-V label (158 modules) & kill if coverage $<$ 50\% or robust SD($d$) $<$ 2.57\,mV & coverage 1.00 (5,688 cells); SD 3.26\,mV & PASS, narrow \\
T37 v1 (tool v0.3; 10 modules) & exclude if residual $>$ 5\% of a cell; kill if $>$ 30\% excluded & 3 / 10 excluded (30\%); one included module visibly misregistered & PASS at the line \\
T37 v1, UVF & same & 9 / 10 over 5\% & fails as built \\
T37 v2 (tool v0.4; post hoc) & frozen 5\% rule plus grid-to-outline $\leq$ 20\% of a cell & 4 / 10 excluded by the 5\% rule; grid-to-outline over 20\% on all 10 & FAIL \\
T02, T12 & & not run & \\
\bottomrule
\end{tabular}
\end{table}

\paragraph{Decision (about 13:30).} The study's lead stopped the A1 arm at T37 as pre-registered (option A); T02 and T12 never ran. The
reasons recorded were the T37 failure, the resistive and possibly inverted label, the row-level dark-IR signal of pick 1, and that a
third, post hoc registration would add a forking path for a likely null. The arm was redirected to a new plan, pick3b on A3, frozen at
13:43 before A3 was extracted: does the EL dark fraction track measured power loss (3b-R1), does drone IR track it and add to EL
(3b-R2), and is there a cell-level IR signal by induced cause (3b-R3, interconnect modules primary). Discovery found no module type in
the release, so under the frozen fallback the primary label became fill-factor (FF) loss; and one EL image per module, so the
two-current label (3b-R4) was not run. No pick3b statistic was computed on 4 November.

\section{4 November 13:51 to 12 November: abrupt end and the record}\label{sec:thermal-end}

The session ended abruptly at about 13:51 on 4 November, cause unknown. Traces: null bytes in the pick3b plan file and in the team
log, an empty file for an attempted version 1.1 of the pick3b plan (its content is lost; version 1 governs), and an empty A3 layout
file. No readout file and no frozen copy was affected. At the end, the report to a collaborator had not been sent and the addendum to
the summary was still a draft. Process faults found along the way were all logged and disclosed: times typed by hand ahead of the
clock; registration self-test criteria set post hoc from synthetic results; a results file overwritten with post hoc fields and the
frozen output restored byte-exact; a tool version run on real data before its review (registration only); one version announced and
withdrawn; one visual-check miss caught by the checker; the false premise about the A1 crop; and a scope slip (all 157 A1 IR modules
instead of the approved 10) caught before any download. On 12 November a session record and handbook chapter were written and
independently verified, and the downloaded data were re-verified. The same evening, asked about thermal, we chose no direct
thermal pick.

\section{15 and 16 November: thermal inside the severity track}\label{sec:thermal-severity}

On 15 November we asked about work on thermal images; it was taken up as part of the severity track, which ran the unread A3
analyses of 4 November, not as a change of direction. Before any A3 statistic existed, a dated amendment to the pick3b plan was
written (15 November, 22:28 to 22:34) and not changed after. It was needed because pick3b version 1 let its thermal readout reuse
pick 1's cells only under an outline rule that pick 1's sixth amendment had replaced, and because pick 1's A3 module table was never
written; both measures were computed by pick 1's frozen definitions. The IR score per module is the spread of the hot spots,
$S_\mathrm{IR}$ = 99th percentile minus median of the array inside a 3\% inset, as the median over frames at 600~W/m$^2$ and above.
Statistics used a module bootstrap of 2,000 draws.

The readouts and their verdicts under the pick3b version 1 bars are in Table~\ref{tab:thermal-3b}. The non-diode numbers are the ones
Section~\ref{sec:deploy-thermal} and Section~\ref{sec:deploy-power} report; the all-module rows are added here. On all 31 non-soiling
modules the EL dark fraction is label-invalid against FF loss (its interval includes 0); our reading is that this is a property of the label: the bypass-diode modules have
the most EL darkness and a large maximum-power loss, but FF loss does not register it. The per-cause medians of $S_\mathrm{IR}$ were
0.87\,$^{\circ}$C for cracks, 1.40 for interconnects, 1.13 for potential-induced degradation, 1.57 for glass, 23.9 for bypass diodes,
and 4.2 and 18.3 for the two soiling groups (descriptive). The cell-level readout 3b-R3 was not run: it needs pick 1's sunlit
registration with a co-signed visual check, which was never done, and the track's brief forbade any cell-level thermal claim because the
4 November study had found the dark signal row-level. An independent review on 16 November found no failing item in the track and
reproduced its headline numbers; its wording fixes changed no number.

\begin{table}[htbp]
\centering\small
\caption{Pick3b readouts on A3, computed in the severity track (15 November 2026). Label: FF loss. Spearman $\rho$ with 95\% module
bootstrap interval. Bars of pick3b v1: 3b-R1 label-valid if $\rho \geq 0.5$ with lower bound $> 0$, invalid if $\rho < 0.3$; 3b-R2
``proxy'' with the same bars, ``adds to EL'' if $\Delta R^2 \geq 0.10$ with lower bound $> 0$, ``no added value'' if $\Delta R^2 < 0.05$.}
\label{tab:thermal-3b}
\begin{tabular}{@{}llrll@{}}
\toprule
Readout & Set & $n$ & Value [interval] & Verdict \\
\midrule
3b-R1, EL dark fraction & all non-soiling & 31 & $\rho$ +0.006 [$-$0.380, +0.425] & label-invalid \\
3b-R1, EL dark fraction & non-diode & 24 & $\rho$ +0.636 [+0.236, +0.872] & label-valid \\
3b-R2, $S_\mathrm{IR}$ & all non-soiling & 31 & $\rho$ $-$0.276 [$-$0.575, +0.114] & not a proxy \\
3b-R2, $S_\mathrm{IR}$ & non-diode & 24 & $\rho$ +0.192 [$-$0.305, +0.626] & not a proxy \\
3b-R2, added to EL & all non-soiling & 31 & $\Delta R^2$ $-$0.021 [$-$67.7, +0.248] & no added value \\
3b-R2, added to EL & non-diode & 24 & $\Delta R^2$ +0.102 [$-$0.177, +0.401] & in between \\
\bottomrule
\end{tabular}
\end{table}

\section{What stays open, and what we learned}\label{sec:thermal-open}

Not done as of the last record: pick 1's A3 screening readouts as frozen, the class table with its AUROCs (R1; the severity track gave
only per-cause medians of $S_\mathrm{IR}$) and the recall readout (R2); the sunlit cell-level gate (R3b)
and pick3b's 3b-R3; the A1 extension; the pick 2 visit and the questions to the partner; and the report to a collaborator, whose addendum
remained a draft. Decisions on licences, release and capture funding stayed with a collaborator.

Our reading of the study as a whole:
\begin{enumerate}
\item Thermal is an electrical-consequence signal. It belongs with inactive or disconnected area, hot spots and severity, not with crack
lines or scuff; for scuff against crack, UVF, PL or visible images are better placed.
\item Dark forward-bias IR on the public pairs carried row and module structure, not a cell-level flag (post hoc). The fair tests for a
cell-level inactive flag, sunlit IR on A3 and lock-in IR, were not read.
\item A data set's physics must be known before it is used as a label: A1 is a series-resistance set, and its bright-pixel cell voltage
may invert the direction of ``degraded''.
\item Instruments fail quietly: a passing residual can hide a drifted grid, and an author's crop need not be the module.
\end{enumerate}
